\documentclass[11pt,a4paper]{article}
\usepackage{fontspec,amsmath,amssymb,amsthm,microtype}
\setmainfont{lmroman10-regular.otf}[BoldFont=lmroman10-bold.otf,ItalicFont=lmroman10-italic.otf,BoldItalicFont=lmroman10-bolditalic.otf]
\usepackage[margin=25mm]{geometry}
\usepackage[unicode,hidelinks]{hyperref}
\hypersetup{pdfauthor={LangevinLean Project}}
\setlength{\emergencystretch}{3em}
\setlength{\parskip}{0.4em}
\setlength{\parindent}{0pt}
\allowdisplaybreaks[2]
\raggedbottom
\title{\texorpdfstring{Propagation of finite error moments}{Propagation of finite error moments}}
\author{LangevinLean Project}
\date{}
\begin{document}\maketitle
Draft mathematical exposition; these notes may contain mistakes.


A finite initial coupling error and finite local defects suffice to justify all real transport estimates. Common-noise contraction propagates these finite moments to every iterate.



These estimates justify the Bochner integrals used in the probabilistic local-defect bound. The dominating functions are built from the stated score and noise norms, and the measure-theoretic steps establish measurability together with finite integral, rather than relying on an unproved finite integral.



The key information is about the difference of the coupled states. A domination \(\|A\|\le q\|D\|\) transfers the finite moment of \(D\) to \(A\), and addition of the finite-moment defect then transfers it to the next error. Thus moment propagation does not require separate bounds on every noise term in the error recursion: the shared noise has already canceled. Contractivity \(q<1\) is needed for convergence, but finiteness propagates for every finite nonnegative \(q\).


\section{\texorpdfstring{Finite moments survive a contracted update}{Finite moments survive a contracted update}}


Let \(p\in[0,+\infty]\), \(q\in\mathbb R\), and \(D,A,V,E:\Omega\to H\). Assume \(A\) is almost everywhere strongly measurable, \(E=A-V\) almost everywhere, and \(\|A\|\le q\|D\|\) almost everywhere. If \(D,V\in L^p(P;H)\), then \(E\in L^p(P;H)\).



The domination transfers the finite \(L^p\) moment of \(D\) to \(A\).



The difference \(A-V\) belongs to \(L^p\). Almost everywhere equality then transfers this property to \(E\).



The comparison of norms supplies finite \(p\)-moment size, while almost-everywhere strong measurability supplies the measurability part of membership. Combining both is necessary: a numerical upper bound alone would not define a valid Bochner-integrable random variable. The equality \(E=A-V\) then transfers membership through the algebra of \(L^p\).


\section{\texorpdfstring{All coupling errors have finite second moment}{All coupling errors have finite second moment}}


Assume \(\nabla U\) is \(L\)-Lipschitz, \(q\ge0\), and the initial vector, every noise vector, and every \(Y_k\) are measurable. Suppose

\[Y_{k+1}=T_h(U,h,Y_k,\xi_k)+V_k\quad\text{almost everywhere},\]

and common-noise updates satisfy \(\|T_h(U,h,x,z)-T_h(U,h,y,z)\|\le q\|x-y\|\). If \(X_0-Y_0\in L^2\) and every \(V_k\in L^2\), then \(X_n-Y_n\in L^2\) for all \(n\). This assertion does not require stationarity or \(q<1\).



Name the actual iterates \(X_k\), whose measurability follows from the chain lemma. We must prove \(X_n-Y_n\in L^2\).



Induct on \(n\). The initial case is the assumed finite initial error; for the successor, assume \(X_n-Y_n\in L^2\).



Set \(A=T_h(U,h,X_n,\xi_n)-T_h(U,h,Y_n,\xi_n)\), the update difference before adding the reference defect.



The gradient and both update inputs are measurable, so the update difference \(A\) is measurable and hence almost everywhere strongly measurable.



Subtract the recursive updates on the full-measure set where the reference equation holds. This gives \(X_{n+1}-Y_{n+1}=A-V_n\).



The common-noise contraction gives \(\|A\|\le q\|X_n-Y_n\|\) pointwise, hence almost everywhere.



Apply the finite-moment lemma to the induction hypothesis and \(V_n\in L^2\), proving \(X_{n+1}-Y_{n+1}\in L^2\).



At step \(n\), the shared-noise update difference is dominated by \(q\|X_n-Y_n\|\). The assumed defect has a finite second moment, so their difference does as well. Induction repeats this argument for every finite index. The proof establishes moments even in regimes with \(q\ge1\); the strict upper bound on \(q\) is used only later for decay.


\end{document}
